\subsection{Matrix coefficients}

DEFINITION 5. --- Let G be a topological group and let \(\varrho\) be a unitary representation of G in a Hilbert space E. Let \(x\) and \(y\) be elements of E. The function from G to K given by \(g \mapsto \langle x | \varrho(g)y \rangle\) is called a matrix coefficient of \(\varrho\) , or a representative function. If \(x = y\) , one says that it is a diagonal matrix coefficient. If \(\varrho\) is of finite dimension, we say that it is a finite-dimensional matrix coefficient.

The matrix coefficients of \(\varrho\) are continuous and bounded functions on G. We denote by \(\Upsilon(G)\) (resp. \(\Theta(G)\) ) the set of matrix coefficients of complex unitary representations (resp. complex finite-dimensional representations) of G.

PROPOSITION 5. — Let G be a topological group. The sets \(\Theta(\mathrm{G})\) and \(\Upsilon(\mathrm{G})\) are unital involutive subalgebras of \(\mathcal{C}_{\mathrm{b}}(\mathrm{G})\) .

The constant function 1 is a matrix coefficient of the trivial representation of G on C. Let \(\varrho\) be a unitary representation of G and let \(x\) and \(y\) be vectors in the space of \(\varrho\) . For every \(\lambda \in \mathbf{C}\) and every \(g \in \mathrm{G}\) , we have \(\lambda \langle x \mid \varrho(g)y \rangle = \langle x \mid \varrho(g)(\lambda y) \rangle\) . Moreover, we have

\[
\langle x \mid \varrho (g) y \rangle = \langle \varrho (g ^ {- 1}) x \mid y \rangle = \overline {{\langle y \mid \varrho (g ^ {- 1}) x \rangle}}
\]

for all \(g \in G\) . Consequently, the sets \(\Theta(G)\) and \(\Upsilon(G)\) are stable under multiplication by scalars and under conjugation.

Let \(\varrho_{1}\) and \(\varrho_{2}\) be unitary representations of G; let \((x_{1}, y_{1})\) be vectors in the space of \(\varrho_{1}\) and \((x_{2}, y_{2})\) be vectors in the space of \(\varrho_{2}\) . For every \(g \in G\) , we then have

\[
\begin{array}{c} \langle x _ {1}   |   \varrho_ {1} (g) y _ {1} \rangle + \langle x _ {2}   |   \varrho_ {2} (g) y _ {2} \rangle = \langle (x _ {1}, x _ {2})   |   (\varrho_ {1} \oplus \varrho_ {2}) (g) (y _ {1}, y _ {2}) \rangle , \\ \langle x _ {1}   |   \varrho_ {1} (g) y _ {1} \rangle \langle x _ {2}   |   \varrho_ {2} (g) y _ {2} \rangle = \langle x _ {1} \otimes x _ {2}   |   (\varrho_ {1} \otimes \varrho_ {2}) (g) (y _ {1} \otimes y _ {2}) \rangle , \end{array}
\]

which proves that \(\Theta(G)\) and \(\Upsilon(G)\) are stable under addition and multiplication. The proposition follows.

